\documentclass[11pt]{article}
\usepackage[margin=1in]{geometry}
\usepackage{amsmath,amssymb,bm}
\usepackage{booktabs}
\usepackage{array}
\usepackage{hyperref}
\usepackage{enumitem}

\title{Mathematical Formulation of the Current AlignNhc Filter}
\author{Derived from the implementation in \texttt{align/src/align\_nhc.rs}}
\date{\today}

\begin{document}
\maketitle

\tableofcontents
\newpage

\section{Purpose}

This document describes the mathematical formulation implemented by the current
\texttt{AlignNhc} filter in this repository. The goal of the filter is to estimate
vehicle attitude, vehicle velocity, IMU-to-vehicle mounting attitude, and IMU
biases using:

\begin{itemize}[leftmargin=1.5em]
\item IMU specific force and angular rate,
\item GNSS velocity in navigation coordinates,
\item a non-holonomic constraint (NHC),
\item a planar-motion gyro constraint.
\end{itemize}

This write-up follows the code as implemented today. It is therefore a
description of the actual estimator, not an idealized estimator.

\section{Frames and Conventions}

The filter uses three frames:

\begin{itemize}[leftmargin=1.5em]
\item \(n\): navigation frame,
\item \(v\): vehicle frame,
\item \(b\): IMU body frame.
\end{itemize}

The two principal attitudes are:

\begin{align}
\mathbf{C}_{n}^{b} &= \text{rotation from body to navigation}, \\
\mathbf{C}_{v}^{b} &= \text{rotation from body to vehicle}.
\end{align}

In the code, these are stored as quaternions:

\begin{align}
\mathbf{q}_{nb} &\leftrightarrow \mathbf{C}_{n}^{b}, \\
\mathbf{q}_{vb} &\leftrightarrow \mathbf{C}_{v}^{b}.
\end{align}

The code uses the convention that \(\mathbf{C}_{a}^{b}\mathbf{x}^{a}\) expresses a vector
from frame \(a\) in frame \(b\). This can be verified from:

\begin{align}
\mathbf{v}^{b} = \mathbf{C}_{n}^{b^\top}\mathbf{v}^{n},
\qquad
\mathbf{v}^{v} = \mathbf{C}_{v}^{b}\mathbf{v}^{b},
\end{align}

which matches the implementation of the NHC prediction.

\section{State Definition}

The nominal state is
\begin{align}
\mathbf{x} =
\begin{bmatrix}
\mathbf{q}_{nb} \\
\mathbf{v}^{n} \\
\mathbf{q}_{vb} \\
\mathbf{b}_{g} \\
\mathbf{b}_{a}
\end{bmatrix},
\end{align}
where:

\begin{itemize}[leftmargin=1.5em]
\item \(\mathbf{q}_{nb}\in \mathbb{H}\): navigation-to-body attitude,
\item \(\mathbf{v}^{n}\in \mathbb{R}^{3}\): vehicle velocity in navigation frame,
\item \(\mathbf{q}_{vb}\in \mathbb{H}\): vehicle-to-body mount attitude,
\item \(\mathbf{b}_{g}\in \mathbb{R}^{3}\): gyro bias in body frame,
\item \(\mathbf{b}_{a}\in \mathbb{R}^{3}\): accelerometer bias in body frame.
\end{itemize}

The linearized error state is 15-dimensional:
\begin{align}
\delta \mathbf{x} =
\begin{bmatrix}
\delta \bm{\theta}_{nb} \\
\delta \mathbf{v}^{n} \\
\delta \bm{\theta}_{vb} \\
\delta \mathbf{b}_{g} \\
\delta \mathbf{b}_{a}
\end{bmatrix}
\in \mathbb{R}^{15}.
\end{align}

The covariance is
\begin{align}
\mathbf{P}\in \mathbb{R}^{15\times 15}.
\end{align}

\section{Initialization}

\subsection{Stationary Initialization Inputs}

The filter initializes from stationary IMU samples:

\begin{align}
\{\mathbf{f}_{m,b}^{(k)}\}_{k=1}^{N}, \qquad
\{\bm{\omega}_{m,b}^{(k)}\}_{k=1}^{N},
\end{align}

and optionally a yaw seed and a reference direction \(\mathbf{x}_{\text{ref}}\).

Define the sample means:
\begin{align}
\bar{\mathbf{f}}_{b} &= \frac{1}{N}\sum_{k=1}^{N}\mathbf{f}_{m,b}^{(k)}, \\
\bar{\bm{\omega}}_{b} &= \frac{1}{N}\sum_{k=1}^{N}\bm{\omega}_{m,b}^{(k)}.
\end{align}

\subsection{Vehicle Down Axis from Gravity}

Under stationarity, the measured specific force is approximately opposite
gravity in body coordinates. The code therefore sets:
\begin{align}
\hat{\mathbf{z}}_{v}^{\,b} = - \frac{\bar{\mathbf{f}}_{b}}{\|\bar{\mathbf{f}}_{b}\|}.
\end{align}

This means the vehicle down axis, expressed in body coordinates, is aligned
with the negative normalized accelerometer mean.

\subsection{Vehicle Forward Axis from a Reference Direction}

Because gravity does not determine yaw, the initialization projects an external
reference direction \(\mathbf{x}_{\text{ref}}\) into the plane orthogonal to
\(\hat{\mathbf{z}}_{v}^{\,b}\):
\begin{align}
\tilde{\mathbf{x}}_{v}^{\,b}
=
\mathbf{x}_{\text{ref}}
-
\hat{\mathbf{z}}_{v}^{\,b}
\left(\hat{\mathbf{z}}_{v}^{\,b}\right)^{\top}
\mathbf{x}_{\text{ref}}.
\end{align}

Then it normalizes this vector and constructs:
\begin{align}
\hat{\mathbf{y}}_{v}^{\,b} &=
\frac{\hat{\mathbf{z}}_{v}^{\,b} \times \hat{\mathbf{x}}_{v}^{\,b}}
{\|\hat{\mathbf{z}}_{v}^{\,b} \times \hat{\mathbf{x}}_{v}^{\,b}\|}, \\
\hat{\mathbf{x}}_{v}^{\,b} &=
\frac{\hat{\mathbf{y}}_{v}^{\,b} \times \hat{\mathbf{z}}_{v}^{\,b}}
{\|\hat{\mathbf{y}}_{v}^{\,b} \times \hat{\mathbf{z}}_{v}^{\,b}\|}.
\end{align}

These three body-expressed vehicle axes form the rows of the tilt-only body to
vehicle rotation:
\begin{align}
\mathbf{C}_{v}^{b,\text{tilt}} =
\begin{bmatrix}
\left(\hat{\mathbf{x}}_{v}^{\,b}\right)^{\top} \\
\left(\hat{\mathbf{y}}_{v}^{\,b}\right)^{\top} \\
\left(\hat{\mathbf{z}}_{v}^{\,b}\right)^{\top}
\end{bmatrix}.
\end{align}

\subsection{Mount-Yaw Preservation}

The code preserves the yaw of an initial mount seed \(\mathbf{q}_{vb,\text{seed}}\).
Let \(\mathbf{C}_{v}^{b,\text{seed}}\) be the corresponding rotation matrix, and
let its yaw be
\begin{align}
\psi_{vb,\text{seed}} = \text{atan2}\left(
\mathbf{C}_{v}^{b,\text{seed}}(2,1),
\mathbf{C}_{v}^{b,\text{seed}}(1,1)
\right).
\end{align}

It then forms a pure vehicle-frame yaw rotation
\begin{align}
\mathbf{C}_{\Delta,vb}
=
\begin{bmatrix}
\cos \psi_{vb,\text{seed}} & -\sin \psi_{vb,\text{seed}} & 0 \\
\sin \psi_{vb,\text{seed}} &  \cos \psi_{vb,\text{seed}} & 0 \\
0 & 0 & 1
\end{bmatrix},
\end{align}
and defines the initialized mount matrix as
\begin{align}
\mathbf{C}_{v}^{b}
=
\mathbf{C}_{v}^{b,\text{tilt}}
\mathbf{C}_{\Delta,vb}.
\end{align}

\subsection{Navigation Attitude from Vehicle Yaw Seed}

The stationary vehicle is assumed level, so the navigation-to-vehicle attitude
is initialized as yaw-only:
\begin{align}
\mathbf{C}_{n}^{v}
=
\begin{bmatrix}
\cos \psi_{nv} & -\sin \psi_{nv} & 0 \\
\sin \psi_{nv} &  \cos \psi_{nv} & 0 \\
0 & 0 & 1
\end{bmatrix},
\end{align}
where \(\psi_{nv}\) is the supplied initial yaw seed.

The navigation-to-body attitude is then
\begin{align}
\mathbf{C}_{n}^{b} = \mathbf{C}_{n}^{v}\mathbf{C}_{v}^{b}.
\end{align}

\subsection{Initial Biases}

The initial gyro bias is set from the stationary gyro average:
\begin{align}
\mathbf{b}_{g,0} = \bar{\bm{\omega}}_{b}.
\end{align}

The initial accelerometer bias is
\begin{align}
\mathbf{b}_{a,0}
=
\bar{\mathbf{f}}_{b}
+
\left(\mathbf{C}_{n}^{b}\right)^{\top}
\begin{bmatrix}
0 \\ 0 \\ g
\end{bmatrix},
\end{align}
where \(g = 9.80665\ \text{m/s}^2\).

\section{Continuous-Time Mechanization}

\subsection{Bias-Corrected IMU Signals}

The code forms:
\begin{align}
\bm{\omega}_{b} &= \bm{\omega}_{m,b} - \mathbf{b}_{g}, \\
\mathbf{f}_{b} &= \mathbf{f}_{m,b} - \mathbf{b}_{a}.
\end{align}

\subsection{Attitude Propagation}

The navigation-to-body quaternion is propagated using a small-angle increment:
\begin{align}
\delta \mathbf{q}(\Delta t)
\approx
\begin{bmatrix}
1 \\
\frac{1}{2}\bm{\omega}_{b}\Delta t
\end{bmatrix},
\end{align}
and the nominal attitude update is
\begin{align}
\mathbf{q}_{nb}^{+}
=
\text{norm}\left(
\mathbf{q}_{nb}\otimes \delta \mathbf{q}(\Delta t)
\right).
\end{align}

\subsection{Velocity Propagation}

The corrected specific force is rotated into the navigation frame:
\begin{align}
\mathbf{a}^{n}
=
\mathbf{C}_{n}^{b}\mathbf{f}_{b}
+
\begin{bmatrix}
0 \\ 0 \\ g
\end{bmatrix}.
\end{align}

Velocity is integrated as
\begin{align}
\mathbf{v}^{n+}
=
\mathbf{v}^{n}
+
\mathbf{a}^{n}\Delta t.
\end{align}

\subsection{Mount Attitude and Bias Propagation}

The nominal mount attitude and the biases are modeled as constants during
prediction:
\begin{align}
\mathbf{q}_{vb}^{+} &= \mathbf{q}_{vb}, \\
\mathbf{b}_{g}^{+} &= \mathbf{b}_{g}, \\
\mathbf{b}_{a}^{+} &= \mathbf{b}_{a}.
\end{align}

\section{Linearized Error-State Dynamics}

The filter uses a first-order linearized error-state model
\begin{align}
\delta \dot{\mathbf{x}} = \mathbf{F}\delta \mathbf{x} + \mathbf{G}\mathbf{w},
\end{align}
discretized as
\begin{align}
\bm{\Phi} \approx \mathbf{I} + \mathbf{F}\Delta t.
\end{align}

\subsection{State Transition Matrix}

Let \([\cdot]_{\times}\) denote the skew-symmetric matrix. Using the current
nominal trajectory, the implemented \(\mathbf{F}\) is:
\begin{align}
\mathbf{F}
=
\begin{bmatrix}
-[\bm{\omega}_{b}]_{\times} & \mathbf{0} & \mathbf{0} & -\mathbf{I} & \mathbf{0} \\
-[\mathbf{a}^{n}]_{\times} & \mathbf{0} & \mathbf{0} & \mathbf{0} & -\mathbf{C}_{n}^{b} \\
\mathbf{0} & \mathbf{0} & \mathbf{0} & \mathbf{0} & \mathbf{0} \\
\mathbf{0} & \mathbf{0} & \mathbf{0} & \mathbf{0} & \mathbf{0} \\
\mathbf{0} & \mathbf{0} & \mathbf{0} & \mathbf{0} & \mathbf{0}
\end{bmatrix}.
\end{align}

Notably:

\begin{itemize}[leftmargin=1.5em]
\item there is no deterministic coupling in the mount-attitude rows,
\item there is no deterministic bias-dynamics model beyond random walk.
\end{itemize}

\subsection{Process Noise Injection}

The code uses a block-diagonal continuous noise covariance:
\begin{align}
\mathbf{Q}_{c}
=
\text{diag}\left(
\sigma_{q,att}^{2}\mathbf{I}_{3},
\sigma_{q,v}^{2}\mathbf{I}_{3},
\sigma_{q,mount}^{2}\mathbf{I}_{3},
\sigma_{q,bg}^{2}\mathbf{I}_{3},
\sigma_{q,ba}^{2}\mathbf{I}_{3}
\right).
\end{align}

The noise input matrix is
\begin{align}
\mathbf{G}
=
\begin{bmatrix}
-\mathbf{I} & \mathbf{0} & \mathbf{0} & \mathbf{0} & \mathbf{0} \\
\mathbf{0} & -\mathbf{C}_{n}^{b} & \mathbf{0} & \mathbf{0} & \mathbf{0} \\
\mathbf{0} & \mathbf{0} & \mathbf{I} & \mathbf{0} & \mathbf{0} \\
\mathbf{0} & \mathbf{0} & \mathbf{0} & \mathbf{I} & \mathbf{0} \\
\mathbf{0} & \mathbf{0} & \mathbf{0} & \mathbf{0} & \mathbf{I}
\end{bmatrix}.
\end{align}

Then the discrete process covariance is
\begin{align}
\mathbf{Q}_{d} = \mathbf{G}\mathbf{Q}_{c}\mathbf{G}^{\top}\Delta t,
\end{align}
and the covariance propagation is
\begin{align}
\mathbf{P}^{+} = \bm{\Phi}\mathbf{P}\bm{\Phi}^{\top} + \mathbf{Q}_{d}.
\end{align}

\section{Measurement Models}

\subsection{GNSS Velocity}

The measurement is GNSS velocity in navigation coordinates:
\begin{align}
\mathbf{z}_{v}^{n} =
\begin{bmatrix}
v_{N} \\ v_{E} \\ v_{D}
\end{bmatrix}.
\end{align}

The prediction is simply the state velocity:
\begin{align}
\hat{\mathbf{z}}_{v}^{n} = \mathbf{v}^{n}.
\end{align}

The Jacobian of each scalar component is
\begin{align}
\mathbf{H}_{v,i}
=
\begin{bmatrix}
\mathbf{0}_{1\times 3} &
\mathbf{e}_{i}^{\top} &
\mathbf{0}_{1\times 9}
\end{bmatrix},
\end{align}
with variance
\begin{align}
R_{v} = \sigma_{gnss}^{2}.
\end{align}

\subsection{Non-Holonomic Constraint}

\subsubsection{Prediction}

The NHC assumes the vehicle has negligible lateral and vertical velocity in the
vehicle frame. First rotate navigation velocity into the body frame:
\begin{align}
\mathbf{v}^{b} = \mathbf{C}_{n}^{b^\top}\mathbf{v}^{n},
\end{align}
then into the vehicle frame:
\begin{align}
\mathbf{v}^{v} = \mathbf{C}_{v}^{b}\mathbf{v}^{b}.
\end{align}

The predicted constrained channels are
\begin{align}
\hat{\mathbf{h}}_{nhc}
=
\begin{bmatrix}
v_{y}^{v} \\
v_{z}^{v}
\end{bmatrix}.
\end{align}

The measurement is implicitly zero:
\begin{align}
\mathbf{z}_{nhc} =
\begin{bmatrix}
0 \\ 0
\end{bmatrix}.
\end{align}

\subsubsection{Jacobian}

Let
\begin{align}
\mathbf{S}_{yz}
=
\begin{bmatrix}
0 & 1 & 0 \\
0 & 0 & 1
\end{bmatrix}.
\end{align}

Then the implemented Jacobian blocks are:
\begin{align}
\mathbf{H}_{\theta nb}
&=
\mathbf{S}_{yz}\mathbf{C}_{v}^{b}[\mathbf{v}^{b}]_{\times}, \\
\mathbf{H}_{v}
&=
\mathbf{S}_{yz}\mathbf{C}_{v}^{b}\mathbf{C}_{n}^{b^\top}, \\
\mathbf{H}_{\theta vb}
&=
\mathbf{S}_{yz}\mathbf{C}_{v}^{b}\left(-[\mathbf{v}^{b}]_{\times}\right).
\end{align}

Therefore the full row Jacobian is
\begin{align}
\mathbf{H}_{nhc}
=
\begin{bmatrix}
\mathbf{H}_{\theta nb} &
\mathbf{H}_{v} &
\mathbf{H}_{\theta vb} &
\mathbf{0}_{2\times 3} &
\mathbf{0}_{2\times 3}
\end{bmatrix}.
\end{align}

\paragraph{Mount yaw scaling in straight motion.}
When the planar-gyro update is not active, the code scales the mount-yaw
column of the NHC Jacobian:
\begin{align}
H_{nhc}(:, 9) \leftarrow \alpha_{straight}\, H_{nhc}(:, 9),
\end{align}
where
\begin{align}
\alpha_{straight} = \texttt{nhc\_straight\_mount\_yaw\_scale}.
\end{align}

In the current default implementation, \(\alpha_{straight}=0.1\), so straight
segments still update mount yaw, but with reduced sensitivity.

\subsection{Planar Gyro Constraint}

\subsubsection{Prediction}

Bias-correct the measured angular rate:
\begin{align}
\bm{\omega}^{b} = \bm{\omega}_{m,b} - \mathbf{b}_{g}.
\end{align}

Rotate into the vehicle frame:
\begin{align}
\bm{\omega}^{v} = \mathbf{C}_{v}^{b}\bm{\omega}^{b}.
\end{align}

Under the planar-motion assumption, roll and pitch rates in the vehicle frame
should be near zero. Thus the predicted constrained channels are
\begin{align}
\hat{\mathbf{h}}_{pg}
=
\begin{bmatrix}
\omega_{x}^{v} \\
\omega_{y}^{v}
\end{bmatrix},
\qquad
\mathbf{z}_{pg} =
\begin{bmatrix}
0 \\ 0
\end{bmatrix}.
\end{align}

\subsubsection{Jacobian}

Let
\begin{align}
\mathbf{S}_{xy}
=
\begin{bmatrix}
1 & 0 & 0 \\
0 & 1 & 0
\end{bmatrix}.
\end{align}

Then the implemented Jacobian blocks are:
\begin{align}
\mathbf{H}_{bg}
&=
\mathbf{S}_{xy}\left(-\mathbf{C}_{v}^{b}\right), \\
\mathbf{H}_{\theta vb}
&=
\mathbf{S}_{xy}\mathbf{C}_{v}^{b}\left(-[\bm{\omega}^{b}]_{\times}\right).
\end{align}

Thus
\begin{align}
\mathbf{H}_{pg}
=
\begin{bmatrix}
\mathbf{0}_{2\times 3} &
\mathbf{0}_{2\times 3} &
\mathbf{H}_{\theta vb} &
\mathbf{H}_{bg} &
\mathbf{0}_{2\times 3}
\end{bmatrix}.
\end{align}

This is important: there is no direct navigation-attitude term in the planar
gyro update, and no direct accelerometer-bias term either.

\section{Measurement Gating}

\subsection{NHC Gate}

The NHC update is enabled if horizontal speed exceeds a threshold:
\begin{align}
\sqrt{v_{N}^{2}+v_{E}^{2}} \ge v_{min,nhc}.
\end{align}

\subsection{Planar Gyro Gate}

The planar gyro update is enabled if all of the following hold:
\begin{align}
\sqrt{v_{N}^{2}+v_{E}^{2}} &\ge v_{min,pg}, \\
|\omega_{z}^{v}| &\ge \omega_{min,pg}, \\
\sqrt{(\omega_{x}^{v})^{2}+(\omega_{y}^{v})^{2}}
&\le
\rho_{max}\, |\omega_{z}^{v}|.
\end{align}

So the planar gyro update is only active during sufficiently planar turns.

\section{Kalman Update}

Each scalar measurement is processed sequentially.

Given innovation \(r\), row Jacobian \(\mathbf{h}\), and variance \(R\),
the code computes:
\begin{align}
\mathbf{p}_{h} &= \mathbf{P}\mathbf{h}^{\top}, \\
S &= R + \mathbf{h}\mathbf{P}\mathbf{h}^{\top}, \\
\mathbf{K} &= \mathbf{p}_{h}/S, \\
\delta \mathbf{x} &= \mathbf{K}r.
\end{align}

The nominal state is corrected by:
\begin{align}
\mathbf{q}_{nb}^{+}
&=
\text{norm}\left(
\mathbf{q}_{nb}\otimes
\delta \mathbf{q}(\delta \bm{\theta}_{nb})
\right), \\
\mathbf{v}^{n+} &= \mathbf{v}^{n} + \delta \mathbf{v}^{n}, \\
\mathbf{q}_{vb}^{+}
&=
\text{norm}\left(
\mathbf{q}_{vb}\otimes
\delta \mathbf{q}(\delta \bm{\theta}_{vb})
\right), \\
\mathbf{b}_{g}^{+} &= \mathbf{b}_{g} + \delta \mathbf{b}_{g}, \\
\mathbf{b}_{a}^{+} &= \mathbf{b}_{a} + \delta \mathbf{b}_{a}.
\end{align}

The covariance update is:
\begin{align}
\mathbf{P}^{+}
=
\mathbf{P} - \mathbf{K}\mathbf{h}\mathbf{P},
\end{align}
followed by explicit symmetrization:
\begin{align}
\mathbf{P}^{+} \leftarrow \frac{1}{2}\left(\mathbf{P}^{+} + \mathbf{P}^{+\top}\right).
\end{align}

\section{Update Sequence}

For each GNSS epoch, the implemented update order is:

\begin{enumerate}[leftmargin=1.8em]
\item integrate all IMU packets since the previous GNSS epoch,
\item apply GNSS velocity update,
\item apply NHC update if the NHC gate passes,
\item apply planar-gyro update if the planar-turn gate passes.
\end{enumerate}

This is exactly how \texttt{update\_all()} is structured.

\section{Yaw Seeding}

After initialization, once horizontal speed exceeds \(5/3.6\ \text{m/s}\), the
filter seeds navigation yaw from GNSS course:
\begin{align}
\psi_{course} = \text{atan2}(v_{E}, v_{N}).
\end{align}

The code sets:
\begin{align}
\mathbf{C}_{n}^{b}
=
\mathbf{C}_{n}^{v}(\psi_{course})\, \mathbf{C}_{v}^{b},
\end{align}
and assigns a yaw variance of \((5^{\circ})^{2}\) to the navigation attitude.

There is also a method to seed mount yaw from course, but it is not used in the
main replay path analyzed here.

\section{Default Tuning}

The current defaults are:

\begin{center}
\begin{tabular}{@{}ll@{}}
\toprule
Parameter & Default \\
\midrule
\(\sigma_{q,att}\) & \(0.01^\circ\) \\
\(\sigma_{q,v}\) & \(0.25\ \text{m/s}\) \\
\(\sigma_{q,mount}\) & \(0.001^\circ\) \\
\(\sigma_{q,bg}\) & \(0.002^\circ/\text{s}\) \\
\(\sigma_{q,ba}\) & \(0.05\ \text{m/s}^{2}\) \\
\(\sigma_{gnss}\) & \(0.3\ \text{m/s}\) \\
\(\sigma_{nhc}\) & \(0.15\ \text{m/s}\) \\
\(\sigma_{pg}\) & \(0.3^\circ/\text{s}\) \\
\(v_{min,nhc}\) & \(3/3.6\ \text{m/s}\) \\
\(v_{min,pg}\) & \(3/3.6\ \text{m/s}\) \\
\(\omega_{min,pg}\) & \(2^\circ/\text{s}\) \\
\(\rho_{max}\) & \(0.25\) \\
\(\alpha_{straight}\) & \(0.1\) \\
\bottomrule
\end{tabular}
\end{center}

\section{Interpretation and Observability}

\subsection{What the Filter Is Actually Estimating}

This is not a mount-only filter. It jointly estimates:

\begin{align}
\mathbf{q}_{nb}, \quad \mathbf{v}^{n}, \quad \mathbf{q}_{vb}, \quad \mathbf{b}_{g}, \quad \mathbf{b}_{a}.
\end{align}

As a result, some residual that one might hope would correct mount attitude can
instead be absorbed by navigation attitude or biases.

\subsection{Where Mount Yaw Comes From}

Mount yaw is only weakly observed in the current formulation.

\begin{itemize}[leftmargin=1.5em]
\item The planar gyro update does \emph{not} directly constrain mount yaw.
It constrains only \(\omega_{x}^{v}\) and \(\omega_{y}^{v}\).
\item Mount yaw is mainly corrected through the NHC velocity measurement.
\item The NHC sensitivity to yaw scales with the body-frame velocity magnitude.
At low speed, yaw observability is weak even in turns.
\item Navigation yaw and mount yaw appear in the NHC Jacobian with opposite-sign
structure, so the two can compete to explain the same residual.
\end{itemize}

\subsection{Covariance Collapse}

Because the mount process noise is very small and because mount yaw is still
updated, though downweighted, during straight motion, the mount covariance can
shrink substantially before late informative maneuvers occur. Once this happens,
later turns may produce residuals but only small state corrections.

\section{Known Limitations of the Current Formulation}

The following properties are direct consequences of the current implementation:

\begin{enumerate}[leftmargin=1.8em]
\item Mount yaw has no direct planar-turn measurement.
\item Straight segments still update mount yaw through the NHC model unless
\(\alpha_{straight}\) is set to zero.
\item Navigation yaw and mount yaw are jointly estimated and not strongly
separated.
\item The same mount process noise is used for roll, pitch, and yaw.
\item The Joseph stabilized covariance form is not used; the implementation uses
\(\mathbf{P}^{+} = \mathbf{P} - \mathbf{K}\mathbf{H}\mathbf{P}\) followed by symmetrization.
\end{enumerate}

\section{Summary}

The implemented \texttt{AlignNhc} filter is a 15-state error-state EKF-like
estimator with:

\begin{itemize}[leftmargin=1.5em]
\item nominal states \((\mathbf{q}_{nb}, \mathbf{v}^{n}, \mathbf{q}_{vb}, \mathbf{b}_{g}, \mathbf{b}_{a})\),
\item GNSS velocity updates,
\item non-holonomic velocity constraints in the vehicle frame,
\item planar-motion gyro constraints in the vehicle frame,
\item stationary gravity-based initialization of mount roll and pitch,
\item GNSS-course-based seeding of navigation yaw.
\end{itemize}

Its strongest practical behaviors are:

\begin{itemize}[leftmargin=1.5em]
\item mount roll and pitch are fairly directly constrained by gravity and NHC,
\item mount yaw is only indirectly and weakly constrained,
\item late yaw corrections depend on sufficiently informative turns and on the
filter not having already become overconfident.
\end{itemize}

\end{document}
